\documentclass[11pt, a4paper]{article}

\usepackage[a4paper, top=2.5cm, bottom=2.5cm, left=2cm, right=2cm]{geometry}
\usepackage{amsmath}
\usepackage{amssymb}
\usepackage[colorlinks=true, linkcolor=blue, urlcolor=blue, citecolor=blue]{hyperref}

\title{Theory of Everything - Volume 32 \\ \Large Cybernetics \& Control Theory}
\author{Omega Protocol}
\date{\today}

\begin{document}

\maketitle

\section*{Layman's Summary}
How do you build a robot that can walk without falling over? How does your body keep your temperature perfectly at 98.6 degrees? Volume 32 explores Cybernetics—the science of control. We prove the absolute mathematical limits of how systems can steer themselves through a chaotic universe.

\section{Control Theory (Axiom 35)}
"Ashby's Law of Requisite Variety" states that a control system must be as complex as the environment it is trying to control. We establish this as a core axiom of the universe. An informational network can only survive and regulate its surroundings if its internal mathematical complexity matches the chaos outside. You cannot control what you cannot compute.

\section{Negative Feedback (Theorem)}
When you get hot, you sweat to cool down. This is negative feedback. We prove that all forms of "homeostasis" (stable balance) are achieved strictly through topological loops that invert the sign of an incoming informational signal. The universe maintains stability by mathematically fighting back against changes.

\section{Feedforward Anticipation (Theorem)}
The most advanced networks don't just react; they predict. We prove that complex systems, driven by the "Free Energy Principle," will inevitably build internal, virtual models of their environment. By simulating the future inside their own network, they minimize the mathematical "surprise" of future events, effectively steering reality before it happens.

\end{document}
